\documentclass[11pt,a4paper]{article}

\usepackage[T1]{fontenc}
\usepackage[utf8]{inputenc}
\usepackage[a4paper,margin=27mm]{geometry}
\usepackage{amsmath,amssymb,amsthm,mathtools}
\usepackage{newpxtext,newpxmath}
\usepackage{microtype}
\usepackage{booktabs,tabularx,array}
\usepackage{enumitem}
\usepackage{xcolor}
\usepackage{titlesec}
\usepackage{mdframed}
\usepackage[authoryear,round]{natbib}
\usepackage[colorlinks=true,linkcolor=blue!45!black,
  citecolor=red!45!black,urlcolor=blue!45!black]{hyperref}

\definecolor{otblue}{HTML}{1F4E79}
\definecolor{otgray}{HTML}{F1F3F5}
\definecolor{otline}{HTML}{AAB5BF}
\titleformat{\section}{\normalfont\Large\bfseries\color{otblue}}
  {\thesection}{.7em}{}
\titleformat{\subsection}{\normalfont\large\bfseries\color{otblue!80!black}}
  {\thesubsection}{.65em}{}
\titleformat{\paragraph}[runin]
  {\normalfont\bfseries\color{otblue!75!black}}{}{0pt}{}
\titlespacing*{\section}{0pt}{2.7ex plus .7ex minus .2ex}{1.1ex}
\titlespacing*{\subsection}{0pt}{2ex plus .5ex minus .2ex}{.75ex}
\titlespacing*{\paragraph}{0pt}{1.3ex plus .3ex minus .1ex}{.55em}

\newtheorem{theorem}{Theorem}[section]
\newtheorem{proposition}[theorem]{Proposition}
\newtheorem{corollary}[theorem]{Corollary}
\theoremstyle{definition}
\newtheorem{example}[theorem]{Example}
\newtheorem{remark}[theorem]{Remark}
\newmdenv[
  backgroundcolor=otgray,linecolor=otline,linewidth=.5pt,roundcorner=2pt,
  innerleftmargin=8pt,innerrightmargin=8pt,innertopmargin=6pt,
  innerbottommargin=6pt,skipabove=8pt,skipbelow=8pt
]{definitionbox}
\newenvironment{definition}[1][]{%
  \begin{definitionbox}\refstepcounter{theorem}%
  \noindent\textbf{Definition~\thetheorem}%
  \if\relax\detokenize{#1}\relax.\else\ \textbf{(#1).}\fi\enspace
}{\end{definitionbox}}

\newcommand{\RR}{\mathbb R}
\newcommand{\Pp}{\mathcal P}
\newcommand{\Ptwo}{\mathcal P_2}
\newcommand{\Wass}{\mathcal W}
\newcommand{\argmin}{\operatorname*{arg\,min}}
\newcommand{\diverg}{\operatorname{div}}
\newcommand{\dd}{\,\mathrm d}
\newcommand{\al}{\alpha}
\newcommand{\norm}[1]{\left\lVert#1\right\rVert}
\newcommand{\abs}[1]{\left\lvert#1\right\rvert}
\newcommand{\ind}{\iota}

\title{Why Variational Mean Field Games Are Convex\\
  \large Momentum variables, perspective functions, and the Benamou--Brenier structure}
\author{Gabriel Peyr\'e}
\date{July 2026}
\hypersetup{
  pdftitle={Why Variational Mean Field Games Are Convex},
  pdfauthor={Gabriel Peyre},
  pdfsubject={Variational mean field games and convex momentum formulations},
  pdfkeywords={mean field games, optimal transport, Benamou-Brenier, perspective function, convex optimization}
}

\begin{document}
\maketitle

\begin{abstract}
A basic variational mean field game minimizes kinetic effort and congestion
over a population curve whose initial distribution is prescribed and whose
final distribution is penalized. In density and velocity variables, neither
the kinetic energy \(\rho|v|^2/2\) nor the constraint
\(\partial_t\rho+\operatorname{div}(\rho v)=0\) is jointly convex. Convexity
appears after introducing the mass flux \(m=\rho v\): the constraint becomes
affine and the kinetic energy becomes \(|m|^2/(2\rho)\), the perspective of a
quadratic function. This note proves that mechanism carefully, including the
extension at vacuum, the hypotheses on the congestion potential, the convex
dual, and the Hamilton--Jacobi/continuity optimality system. It also separates
this full-horizon problem from Benamou--Brenier transport and JKO minimizing
movements.
\end{abstract}

\tableofcontents

\section{The model and the short answer}

Let \(\mathcal X=\mathbb T^d\), or let \(\mathcal X=\RR^d\) under suitable
moment assumptions. The initial probability measure is
\(\al_0=\rho_0\dd x\). A population curve
\(\al_t=\rho_t\dd x\) is transported by a velocity \(v_t\) according to
\begin{equation}\label{eq:velocity-continuity}
  \partial_t\rho_t+\diverg(\rho_tv_t)=0,
  \qquad \rho_{t=0}=\rho_0.
\end{equation}
The final density is free and is charged through a terminal potential
\(\Psi:\mathcal X\to\RR\). Given a congestion potential
\(G:[0,+\infty)\to\RR\cup\{+\infty\}\), consider
\begin{equation}\label{eq:velocity-primal}
  \inf_{\rho,v}
  \left\{
    \int_0^1\!\int_{\mathcal X}
      \left(\frac12\rho_t\abs{v_t}^2+G(\rho_t)\right)\dd x\dd t
    +\int_{\mathcal X}\Psi\rho_1\dd x
  \right\}
\end{equation}
subject to~\eqref{eq:velocity-continuity} and \(\rho\geq0\). This is a
standard variational MFG model
\citep{BenamouCarlier2015,BenamouCarlierSantambrogio2017}. The factor \(1/2\)
is only a normalization.

\begin{definition}[Momentum or mass flux]
The momentum variable is
\begin{equation}\label{eq:momentum-change}
  m_t(x)=\rho_t(x)v_t(x).
\end{equation}
Thus \(m_t\dd x\) is a vector-valued measure recording transported mass per
unit time, rather than velocity per particle.
\end{definition}

In the variables \((\rho,m)\), the problem becomes
\begin{equation}\label{eq:momentum-primal-preview}
  \inf_{\rho,m}
  \left\{
    \int_0^1\!\int_{\mathcal X}
      \left(\frac{\abs{m_t}^2}{2\rho_t}+G(\rho_t)\right)\dd x\dd t
    +\int_{\mathcal X}\Psi\rho_1\dd x
  \right\},
\end{equation}
subject to the affine conservation law
\begin{equation}\label{eq:momentum-continuity}
  \partial_t\rho_t+\diverg m_t=0,
  \qquad \rho_{t=0}=\rho_0.
\end{equation}
The quotient must be extended correctly at \(\rho=0\), as explained below.

\begin{definition}[Convex variational MFG]
Problem~\eqref{eq:momentum-primal-preview} is a convex optimization problem
provided that \(G\) is convex: the kinetic integrand is jointly convex in
\((\rho,m)\), the congestion is convex in \(\rho\), the terminal cost is
linear, and the constraints are affine.
\end{definition}

\section{Why the velocity formulation is not convex}

The product \(\rho|v|^2\) is convex in each variable separately, but separate
convexity does not imply joint convexity.

\begin{proposition}[Nonconvexity in density--velocity variables]
\label{prop:velocity-nonconvex}
For \(v\neq0\), the map
\[
  (\rho,v)\longmapsto \frac12\rho\abs v^2,\qquad \rho>0,
\]
is not jointly convex. Moreover,
\(\partial_t\rho+\diverg(\rho v)=0\) is nonlinear in \((\rho,v)\).
\end{proposition}

\begin{proof}
In one dimension, the Hessian of \(K(\rho,v)=\rho v^2/2\) is
\[
  \nabla^2K(\rho,v)=
  \begin{pmatrix}0&v\\v&\rho\end{pmatrix},
\]
whose determinant is \(-v^2<0\). It is therefore indefinite whenever
\(v\neq0\). The multidimensional claim follows by restriction to a line.
The product \(\rho v\) in the continuity equation is bilinear.
\end{proof}

The nonlinear change \((\rho,v)\mapsto(\rho,m=\rho v)\) does two things at
once: it linearizes conservation and changes the kinetic integrand. At vacuum,
\(v\) is physically irrelevant, while finite action forces \(m=0\).

\section{The perspective mechanism}\label{sec-perspective}

\begin{definition}[Closed perspective]
Let \(L:\RR^d\to\RR\cup\{+\infty\}\) be proper, convex, and lower
semicontinuous. Its closed perspective is
\begin{equation}\label{eq:general-perspective}
  J_L(a,m)=
  \begin{cases}
    aL(m/a),&a>0,\\
    L^\infty(m),&a=0,\\
    +\infty,&a<0,
  \end{cases}
\end{equation}
where \(L^\infty\) is the recession function of \(L\).
\end{definition}

For the quadratic Lagrangian \(L(v)=|v|^2/2\), this is
\begin{equation}\label{eq:quadratic-perspective}
  J(a,m)=
  \begin{cases}
    \dfrac{\abs m^2}{2a},&a>0,\\[1mm]
    0,&(a,m)=(0,0),\\
    +\infty,&a=0,\ m\neq0,\\
    +\infty,&a<0.
  \end{cases}
\end{equation}
This boundary convention is essential; the symbol \(0/0\) is not a
definition.

\begin{proposition}[Convexity of the perspective]
\label{prop:perspective-convexity}
If \(L\) is proper, convex, and lower semicontinuous, then \(J_L\) is jointly
convex, lower semicontinuous, and positively one-homogeneous:
\[
  J_L(\lambda a,\lambda m)=\lambda J_L(a,m)
  \qquad(\lambda\geq0).
\]
\end{proposition}

\begin{proof}
Take \(a_0,a_1>0\), \(m_0,m_1\in\RR^d\), and \(\theta\in[0,1]\). Set
\[
  a_\theta=(1-\theta)a_0+\theta a_1,\quad
  \lambda_0=\frac{(1-\theta)a_0}{a_\theta},\quad
  \lambda_1=\frac{\theta a_1}{a_\theta}.
\]
Then \(\lambda_0+\lambda_1=1\) and
\[
  \frac{(1-\theta)m_0+\theta m_1}{a_\theta}
  =\lambda_0\frac{m_0}{a_0}+\lambda_1\frac{m_1}{a_1}.
\]
Convexity of \(L\), multiplied by \(a_\theta\), proves the convexity
inequality. The recession value at \(a=0\) is precisely the
lower-semicontinuous closure. Homogeneity follows from the definition.
\end{proof}

For the quadratic case, the epigraph is the rotated second-order cone
\[
  \operatorname{epi}J
  =\{(a,m,s):a\geq0,\ s\geq0,\ \abs m^2\leq2as\},
\]
which gives another direct proof of convexity.

\begin{remark}[A recurring convexity mechanism]
The same perspective transform makes
\(b\varphi(a/b)\) jointly convex in the theory of
\(\varphi\)-divergences. Here it turns a per-particle velocity cost into a
convex density--flux cost.
\end{remark}

\section{The convex primal problem}

Define
\begin{equation}\label{eq:convex-energy}
  \mathcal E(\rho,m)=
  \int_0^1\!\int_{\mathcal X}
    \bigl(J(\rho_t,m_t)+G(\rho_t)\bigr)\dd x\dd t
  +\int_{\mathcal X}\Psi\rho_1\dd x.
\end{equation}
The variational MFG is
\begin{equation}\label{eq:convex-primal}
  \inf\left\{
    \mathcal E(\rho,m):
    \partial_t\rho+\diverg m=0,\ \rho_{t=0}=\rho_0
  \right\}.
\end{equation}

\begin{theorem}[Convexity of the variational MFG]
\label{thm:variational-mfg-convexity}
If \(G:[0,+\infty)\to\RR\cup\{+\infty\}\) is proper, convex, and lower
semicontinuous, then~\eqref{eq:convex-primal} is a convex optimization problem
in \((\rho,m)\).
\end{theorem}

\begin{proof}
The perspective \(J\) is jointly convex by
Proposition~\ref{prop:perspective-convexity}; \(G\) is convex by assumption;
integration preserves convexity; and the endpoint term is linear. The
continuity equation and initial condition are affine constraints.
\end{proof}

If \(G\) is nonconvex, the momentum substitution still convexifies the kinetic
term and linearizes the PDE, but it cannot convexify \(G\). A convex terminal
functional \(\mathcal G(\rho_1\dd x)\) may replace the linear terminal term
without changing the theorem. A fixed endpoint is also a convex affine
constraint.

\begin{proposition}[Basic existence]
\label{prop:variational-mfg-existence}
Let \(\mathcal X=\mathbb T^d\), let \(\Psi\) be lower semicontinuous and
bounded below, and assume
\(G\) is convex, lower semicontinuous, and superlinear:
\[
  G(r)/r\longrightarrow+\infty\qquad(r\to+\infty).
\]
If the infimum is finite, then the lower-semicontinuous relaxation of
\eqref{eq:convex-primal} to density--flux measures admits a minimizer.
\end{proposition}

\begin{proof}
For a minimizing sequence, superlinearity gives uniform integrability of the
densities. The kinetic action controls the Wasserstein metric variation of the
curves. The continuity equation is closed under weak convergence of
density--flux pairs. Convexity, lower semicontinuity, and one-homogeneity of
\(J\) give lower semicontinuity of its integral, while the congestion term is
lower semicontinuous by convexity and uniform integrability. Since \(\Psi\) is
lower semicontinuous and bounded below, the terminal term is lower
semicontinuous under narrow convergence. The direct method produces a
minimizer. On \(\RR^d\), an additional moment or coercivity bound is needed to
prevent escape to infinity.
\end{proof}

\paragraph{Strict convexity.}
The perspective \(J\) is one-homogeneous and hence not strictly convex along
rays. If \(G\) is strictly convex, any two minimizers have the same density.
Once \(\rho>0\) is fixed, strict convexity of
\(m\mapsto|m|^2/(2\rho)\) identifies the flux.

\section{Optimality conditions and the MFG PDE}

The congestion potential \(G\) is not itself the pointwise cost perceived by
one infinitesimal agent. The coupling appearing in the HJB equation is
\[
  g(r)=G'(r)
\]
when \(G\) is differentiable, or a subgradient \(g\in\partial G(r)\) in the
nonsmooth case.

\begin{theorem}[Euler--Lagrange system]
\label{thm:variational-mfg-kkt}
Assume a minimizer \((\rho,m)\) is smooth, \(\rho>0\), and \(G\) is
differentiable. Then there is a potential \(u\) such that
\begin{equation}\label{eq:variational-mfg-system}
  \left\{
  \begin{aligned}
    -\partial_tu+\frac12\abs{\nabla u}^2&=G'(\rho),
      &&u_{t=1}=\Psi,\\
    \partial_t\rho-\diverg(\rho\nabla u)&=0,
      &&\rho_{t=0}=\rho_0,\\
    m&=-\rho\nabla u.
  \end{aligned}
  \right.
\end{equation}
This is the first-order quadratic MFG system with local coupling
\(g(\rho)=G'(\rho)\).
\end{theorem}

\begin{proof}
Introduce the multiplier \(-u\) for the continuity equation. Integration by
parts gives the Lagrangian
\begin{align*}
  \mathscr L(\rho,m,u)
  &=
  \int_0^1\!\int_{\mathcal X}
  \left(
    \frac{\abs m^2}{2\rho}+G(\rho)
    +\rho\partial_tu+m\cdot\nabla u
  \right)\dd x\dd t\\
  &\quad
  +\int_{\mathcal X}(\Psi-u_1)\rho_1\dd x
  +\int_{\mathcal X}u_0\rho_0\dd x.
\end{align*}
Stationarity in \(m\) gives \(m/\rho+\nabla u=0\), hence
\(m=-\rho\nabla u\). Stationarity in \(\rho\) gives
\[
  -\frac{\abs m^2}{2\rho^2}+G'(\rho)+\partial_tu=0.
\]
Substituting the flux relation yields the HJB equation. Since the terminal
density is free and positive, terminal stationarity gives \(u_1=\Psi\).
Primal feasibility gives the continuity equation.
\end{proof}

\begin{remark}[Vacuum and terminal complementarity]
Without positivity, the terminal condition is
\[
  u_1\leq\Psi,\qquad(\Psi-u_1)\rho_1=0.
\]
Thus equality holds on the support of \(\rho_1\). Likewise, \(G'(\rho)\) is
replaced by a subgradient when \(G\) is nonsmooth.
\end{remark}

\paragraph{General kinetic costs.}
For a convex Lagrangian \(L(v)\), use
\(J_L(\rho,m)=\rho L(m/\rho)\). With
\[
  H(p)=\sup_{v\in\RR^d}\{-v\cdot p-L(v)\},
\]
the optimality system becomes
\begin{equation}\label{eq:general-mfg-system}
  -\partial_tu+H(\nabla u)=G'(\rho),
  \qquad
  \partial_t\rho+\diverg\left(-\rho\nabla_pH(\nabla u)\right)=0.
\end{equation}
The quadratic case has \(H(p)=|p|^2/2\).

\section{The convex dual}

Extend \(G\) by \(+\infty\) on negative arguments and define
\[
  G^*(s)=\sup_{r\geq0}\{rs-G(r)\}.
\]
The restriction \(r\geq0\) implies that \(G^*\) is nondecreasing.

\begin{proposition}[Dual problem]
\label{prop:variational-mfg-dual}
Under a standard qualification condition, the dual of
\eqref{eq:convex-primal} is
\begin{equation}\label{eq:variational-mfg-dual}
  \sup_{u:\,u_1\leq\Psi}
  \left\{
    \int_{\mathcal X}u_0\rho_0\dd x
    -
    \int_0^1\!\int_{\mathcal X}
      G^*\left(-\partial_tu+\frac12\abs{\nabla u}^2\right)\dd x\dd t
  \right\}.
\end{equation}
\end{proposition}

\begin{proof}
Use the integrated Lagrangian from
Theorem~\ref{thm:variational-mfg-kkt}. Pointwise in \((t,x)\),
\[
  \inf_m\left\{\frac{\abs m^2}{2\rho}+m\cdot\nabla u\right\}
  =-\frac12\rho\abs{\nabla u}^2.
\]
The remaining density infimum is
\[
  \inf_{\rho\geq0}
  \left\{G(\rho)+\rho\left(
    \partial_tu-\frac12\abs{\nabla u}^2\right)\right\}
  =
  -G^*\left(-\partial_tu+\frac12\abs{\nabla u}^2\right).
\]
The infimum over the free terminal density is finite exactly when
\(u_1\leq\Psi\). The fixed initial density contributes
\(\int u_0\rho_0\), proving the formula.
\end{proof}

At smooth primal--dual optimality, Fenchel equality gives
\[
  -\partial_tu+\frac12\abs{\nabla u}^2=G'(\rho),
\]
which is the HJB equation. The dual is indeed concave: the argument of
\(G^*\) is convex in \(u\), and \(G^*\) is convex and nondecreasing.

\section{Examples of congestion}

\begin{example}[Power congestion]
For \(q>1\),
\[
  G(r)=\frac{r^q}{q},\qquad g(r)=G'(r)=r^{q-1}.
\]
This is strictly convex and superlinear. It penalizes concentration and gives
the compactness used in Proposition~\ref{prop:variational-mfg-existence}.
\end{example}

\begin{example}[Entropy congestion]
Let \(G(r)=r\log r\) for \(r>0\) and \(G(0)=0\). Then
\(g(r)=1+\log r\). Additive constants in the coupling only shift the value
function, so one often writes \(g(r)=\log r\).
\end{example}

\begin{example}[Hard density cap]
For \(\kappa>0\), define
\[
  G(r)=\ind_{[0,\kappa]}(r)
  =
  \begin{cases}
    0,&0\leq r\leq\kappa,\\
    +\infty,&\text{otherwise}.
  \end{cases}
\]
This is convex but nonsmooth. Its subgradient is a pressure that is inactive
below the cap and active where \(\rho=\kappa\).
\end{example}

\begin{example}[A nonconvex choice]
For \(G(r)=-r^2/2\), the kinetic term remains convex after the momentum
change, but the full problem does not. Multiple local minima and concentration
instabilities can occur. The momentum substitution cannot remove
nonconvexity already present in \(G\).
\end{example}

\section{Benamou--Brenier, JKO, and variational MFG}

\subsection{Benamou--Brenier}

Quadratic dynamic optimal transport is
\begin{equation}\label{eq:bb-comparison}
  \frac12\Wass_2^2(\al_0,\al_1)
  =
  \inf_{\rho,m}
  \left\{
    \int_0^1\!\int J(\rho_t,m_t)\dd x\dd t:
    \partial_t\rho+\diverg m=0,\
    \rho_{t=0}=\rho_0,\
    \rho_{t=1}=\rho_1
  \right\}.
\end{equation}
The variational MFG adds the running potential
\(\int_0^1\int G(\rho_t)\) and leaves the endpoint free, selecting it through
\(\int\Psi\rho_1\). It is thus a convex dynamic transport problem, but not a
distance between two prescribed measures. Setting \(G=0\) and fixing the
endpoint recovers Benamou--Brenier \citep{BenamouBrenier2000}.

\subsection{JKO}

A JKO step for an energy \(f\) is
\begin{equation}\label{eq:jko-comparison}
  \al^{k+1}\in\argmin_{\al}
  \left\{
    \frac1{2\tau}\Wass_2^2(\al^k,\al)+f(\al)
  \right\}.
\end{equation}
Using~\eqref{eq:bb-comparison}, it is a one-period planning problem with a
free endpoint and terminal functional \(f(\al_1)\). It has no running
congestion unless one adds it explicitly.

A variational MFG minimizes one full-horizon action containing
\(\int_0^1G(\rho_t)\dd t\) and encodes anticipation over that horizon. A JKO
scheme repeatedly solves short proximal problems; its small-step limit is a
dissipative Wasserstein gradient flow
\citep{JordanKinderlehrerOtto1998,AmbrosioGigliSavare2008}.

\begin{remark}[The common core]
Benamou--Brenier is endpoint-to-endpoint transport, JKO is repeated
short-horizon descent, and a potential MFG is a full-horizon variational game.
All three share the convex mass--momentum perspective.
\end{remark}

\section{Diffusion and discretization}

\paragraph{Diffusion.}
With idiosyncratic Brownian noise of variance \(2\nu\), use
\[
  \partial_t\rho+\diverg m=\nu\Delta\rho.
\]
This is still affine in \((\rho,m)\), so convexity is unchanged. The
optimality system becomes
\[
  -\partial_tu-\nu\Delta u+\frac12\abs{\nabla u}^2=G'(\rho),
  \qquad
  \partial_t\rho-\nu\Delta\rho-\diverg(\rho\nabla u)=0.
\]
For quadratic drift costs, the problem also has an entropy-minimization
interpretation on path space and admits Sinkhorn-like algorithms
\citep{BenamouCarlierDiMarinoNenna2019}.

\paragraph{Discrete convex program.}
After conservative discretization, collect cell densities in \(r\) and edge
fluxes in \(M\). The continuity equation becomes the affine system
\[
  Ar+BM=b.
\]
The objective is a sum of local convex perspective and congestion terms plus
a terminal linear term. Hence the discretized problem is finite-dimensional
and convex. Primal--dual splitting, ADMM, and augmented-Lagrangian methods
exploit the split between local proximal operations and the global linear
conservation constraint \citep{BenamouCarlier2015,AchdouLauriere2020}.

\section{Summary}

The precise answer is
\[
  \boxed{\text{the momentum formulation is convex whenever \(G\) is convex.}}
\]
The reason is not that the original density--velocity problem was secretly
jointly convex. Rather:
\begin{enumerate}[leftmargin=2em]
  \item \(m=\rho v\) makes conservation affine;
  \item the kinetic term becomes the closed perspective
  \(J(\rho,m)=|m|^2/(2\rho)\);
  \item this perspective is jointly convex and lower semicontinuous, including
  at vacuum;
  \item convex congestion and convex endpoint penalties preserve convexity.
\end{enumerate}
If \(G\) is nonconvex, the full problem remains nonconvex. When \(G\) is
convex, the KKT equations are exactly the Hamilton--Jacobi/continuity system
of a potential MFG, with coupling \(g=G'\).

\bibliographystyle{plainnat}
\bibliography{references}

\end{document}
